\section{Data}\label{sec:data}

All data are public, and the code that downloads and assembles them is available in the
replication package. The country-year panel covers the 159 countries with a population of at
least one million in 2010, from 1998 to 2023. Table~\ref{tab:desc} reports descriptive statistics
for the estimation sample.

\paragraph{Decile income growth.} Distributional outcomes come from the World Bank's Poverty and
Inequality Platform (PIP), which reports the mean of the welfare aggregate and the share of each
decile for every household survey in its database. We take national surveys, keep one welfare
concept per country (income or consumption, whichever has more surveys), and form \emph{spells}
from consecutive surveys of the same country that PIP flags as comparable and that are at most ten
years apart. For each spell and decile, the outcome $Y_{idj}$ is the annualised log growth of the
decile's mean welfare, in per cent, where the decile mean equals the survey mean times ten times
the decile share in purchasing-power-parity dollars per day. This yields 1,723 spells in 130
countries. The estimation sample, which requires every covariate to be observed at the start of the
spell and capability to be measured (from 2001), contains \NSpells{} spells in \NCountriesSp{}
countries. Outcomes are winsorised at the 1st and 99th percentiles within each decile.

\paragraph{Temperature shocks.} Annual mean surface temperature for each country is the average of
the monthly series from the ERA5 reanalysis of the Copernicus Climate Change Service, as
aggregated to countries by Our World in Data. Following the climate-econometrics literature
\citep{dell2012,burke2015}, the shock $S_{it}$ is the deviation of annual temperature from a
country-specific linear trend estimated over 1980--2023. Its standard deviation across country-years
is 0.40$^{\circ}$C. For a spell, $S_{ij}$ is the mean of the annual shocks over the years the spell
covers, and $\bar s_i$, the typical shock size used in the welfare aggregation, is the standard
deviation of the country's shocks over 2000--2023. The baseline climate is the mean temperature of
1980--2000.

\paragraph{Climate-AI capability and knowledge sovereignty.} We query the OpenAlex index of
scholarly works \citep{priem2022} for every country and year, counting works with at least one
author affiliated in the country. AI works are those tagged with the ``artificial intelligence'' or
``machine learning'' concepts. \emph{Climate-AI} works are AI works also tagged with ``climate
change''. Because capability accumulates, we construct a knowledge stock
$K_{it}=\sum_{\tau\le t}(1-\delta)^{t-\tau}C_{i\tau}$ from annual climate-AI output $C_{i\tau}$ since
1998, with depreciation $\delta=0.15$, and express it per million population. A country has
climate-AI capability ($T_{it}=1$) if its lagged stock per capita is above the cross-country median
of that year. A spell takes the value of $T$ at its start year. \emph{Knowledge sovereignty} is the
share of a country's climate-AI works in which every author is affiliated in that country, computed
over 1998--2007, 2008--2015 and 2016--2023. It is set to a separate value when a country has fewer
than five climate-AI works in the period. General AI output per million population enters the
covariates, so the capability indicator captures specialisation in climate-relevant AI rather than
AI research in general. Bibliometric output measures research capacity, not deployment. We discuss
what this implies for interpretation in Section~\ref{sec:discussion}.

\paragraph{Covariates.} From the World Development Indicators we take GDP per capita at
purchasing-power parity (log), the Gini index (interpolated linearly within country between
surveys, without extrapolation), the share of the population using the internet, the agricultural
share of value added, the urban population share and population (log). Adaptation readiness, which
summarises economic, governance and social conditions for adaptation investment, comes from the
ND-GAIN index \citep{chen2015ndgain}. For the cost analysis we use per-capita electricity
consumption, per-capita CO$_2$ emissions and the renewable share of electricity generation from Our
World in Data. Countries are classified as emerging and developing economies (EMDEs) if the World
Bank does not classify them as high income.
